\section{Related Work}
\label{sec:related_work}

\paragraph{Text-based document RAG.} Standard RAG chunks extracted text and retrieves with a dense bi-encoder~\citep{dpr2020}. On visually rich documents this requires a layout-aware parser or OCR, both lossy on complex layouts, and discards 2D structure before retrieval ever runs.

\paragraph{Screenshot-based visual RAG.} PixelRAG~\citep{pixelrag2026} retrieves over rendered-page image tiles with a VLM embedding model, part of an established line retrieving over page images rather than text: ColPali~\citep{colpali2024} (late-interaction patch embeddings), Document Screenshot Embedding~\citep{dse2024} (whole-page VLM embeddings, no parsing), and VisRAG~\citep{visrag2024}/M3DocRAG~\citep{m3docrag2024} (retrieve-then-generate pipelines, the latter multi-page). None address how a page is cut into retrievable units for a VLM -- what our tiling (\S\ref{sec:tiling}) targets, so it could feed any of them, not only PixelRAG, whose own blind fixed-height grid lets a table or paragraph straddle a tile boundary. PixelRAG also reports up to $3\times$ token-cost reduction at lower resolution without specifying the filter or resolutions tested; we build a matched full-/reduced-resolution pair on our own tiles to measure this directly (\S\ref{sec:contrastive}).

\paragraph{DOM-based and block-level web page segmentation.} Cutting a page into coherent units predates VLM retrieval: VIPS~\citep{vips2003} segments a page top-down, combining DOM structure with visual cues into a tree of content blocks, and block-based web search~\citep{blockwebsearch2004} shows indexing and query expansion at this block level improves web retrieval over the whole page. Our tiling (\S\ref{sec:tiling}) targets the same problem with a different mechanism -- boundaries come directly from each element's own bounding box plus merge/patch rules, not VIPS's recursive coherence heuristics -- and for a different consumer: a multimodal embedding rather than block-level lexical features.

\paragraph{Graph-structured evidence and controllers.} MAGE-RAG~\citep{magerag2026} (evaluated on LongDocURL~\citep{longdocurl2024}, MMLongBench-Doc~\citep{mmlongbenchdoc2024}) builds an offline multigranular evidence graph and reads it with a budgeted controller -- the direct ancestor of \S\ref{sec:controller}. The difference is provenance: we build the graph from a rendered page's DOM, giving exact boundaries and reading order for free, and instantiate all five of its relation types (\S\ref{sec:graph}). GraphRAG~\citep{graphrag2024} shares the premise of building an explicit graph before querying it, over extracted entities rather than layout structure. Layout- and markup-aware pretraining (LayoutLM~\citep{layoutlm2020}, MarkupLM~\citep{markuplm2022}) instead bakes structure into the model's representation; we use the DOM directly at inference time, with no pretraining required. Web-agent work (WebArena~\citep{webarena2023}, Mind2Web~\citep{mind2web2023}, SeeAct~\citep{seeact2024}) has already studied when DOM structure is reliable enough to act on -- the same failure mode (JS-hydrated, div-soup markup) that scopes our own claim to semantically marked-up pages (\S\ref{sec:limitations}).

\paragraph{Contrastive fine-tuning and hard-negative mining.} Training with in-batch negatives plus mined hard negatives is standard in dense retrieval, e.g.\ DPR's BM25-mined negatives~\citep{dpr2020}; we use TF-IDF at tile granularity instead. Training pairs are generated synthetically by prompting a VLM per tile, in the spirit of InPars~\citep{inpars2022} and Promptagator~\citep{promptagator2022} adapted to image tiles. We fine-tune with LoRA~\citep{lora2021} on both the decoder and the vision tower.
